% ============================================================
%   导言区 / 序言（preamble）—— 管整体风格
%   由 main.tex 用  % ============================================================
%   导言区 / 序言（preamble）—— 管整体风格
%   由 main.tex 用  % ============================================================
%   导言区 / 序言（preamble）—— 管整体风格
%   由 main.tex 用  % ============================================================
%   导言区 / 序言（preamble）—— 管整体风格
%   由 main.tex 用  \input{preamble}  引入（在 \begin{document} 之前）。
%   统一改风格（字体、配色、页面、框）都是改这里。
% ============================================================

% ---------- 字体（XeLaTeX，已由 ctex 加载 xeCJK/fontspec） ----------
% 正文思源宋体；中文「斜体」按国内惯例映射成楷体（\textit / \emph 即楷体）
\setCJKmainfont[ItalicFont=LXGW WenKai, BoldItalicFont=LXGW WenKai]{Noto Serif CJK SC}
\setCJKsansfont{Noto Sans CJK SC}    % 标题中文：思源黑体
\setCJKmonofont{LXGW WenKai Mono}    % 中文等宽：霞鹜文楷 Mono
\setmonofont{JetBrainsMono NF}       % 拉丁等宽（代码）：JetBrains Mono Nerd Font
% 拉丁正文与数学保留 Computer Modern。

% ---------- 数学 ----------
\usepackage{amsmath,amssymb}
\usepackage{bm}
\usepackage{mathtools}
\numberwithin{equation}{chapter}   % 公式编号按章：1.1 / 1.2 ...
\numberwithin{figure}{chapter}     % 图编号按章

% ---------- 页面（A4） ----------
\usepackage[a4paper, top=2.3cm, bottom=2.3cm, left=2.2cm, right=2.2cm,
  headheight=14pt, headsep=6mm, footskip=11mm]{geometry}
\usepackage{setspace}
\onehalfspacing

% ---------- 图形 ----------
\usepackage{graphicx}
\usepackage{tikz}
\usetikzlibrary{calc, arrows.meta, positioning, shapes.geometric}

% ---------- 表格 ----------
\usepackage{booktabs}
\usepackage{array}

% ---------- 物理单位 ----------
\usepackage{siunitx}

% ---------- 列表间距 ----------
\usepackage{enumitem}
\setlist{topsep=3pt, partopsep=0pt, itemsep=2pt, parsep=0pt}

% ---------- 颜色（知识块调色板 + 封面） ----------
\usepackage{xcolor}
\definecolor{bl}{HTML}{3B6BB0}   % 蓝 — 定义
\definecolor{tl}{HTML}{1F9B8E}   % 青 — 理解
\definecolor{gr}{HTML}{4C9F70}   % 绿 — 例
\definecolor{og}{HTML}{E8933A}   % 橙 — 证明思路
\definecolor{pu}{HTML}{7B5EA7}   % 紫 — 命题
\definecolor{rd}{HTML}{C0504D}   % 红 — 定理
\definecolor{in}{HTML}{283593}   % 靛 — 引理
\definecolor{pk}{HTML}{C86B98}   % 粉 — 推论
\definecolor{gy}{HTML}{6B7280}   % 灰 — 注
\definecolor{coverblue}{HTML}{283593}   % 封面深蓝
\definecolor{coverlight}{HTML}{5C6BC0}  % 封面浅蓝
\definecolor{covergray}{HTML}{5A5F6E}   % 封面灰
\definecolor{link}{HTML}{1F6F8B}

% ---------- 页眉页脚 ----------
\usepackage{fancyhdr}
\pagestyle{fancy}
\fancyhf{}
\fancyhead[L]{\small\nouppercase{\rightmark}}
\fancyhead[R]{\small\thepage}
\renewcommand{\headrulewidth}{0.4pt}

% ---------- 章节标题（加大醒目） ----------
\usepackage{titlesec}
\titleformat{\chapter}{\centering\fontsize{20}{26}\selectfont\bfseries}{\Huge 第\chinese{chapter}章}{0.6em}{}
\titlespacing*{\chapter}{0pt}{-16pt}{30pt}
\titleformat{\section}{\fontsize{16}{20}\selectfont\bfseries}{\thesection}{0.6em}{}
\titlespacing*{\section}{0pt}{16pt}{6pt}
\titleformat{\subsection}{\fontsize{14}{18}\selectfont\bfseries}{\thesubsection}{0.6em}{}
\titlespacing*{\subsection}{0pt}{12pt}{4pt}
\titleformat{\subsubsection}{\fontsize{12}{15}\selectfont\bfseries}{\thesubsubsection}{0.6em}{}

% ---------- 目录（加大字号） ----------
\usepackage{tocloft}
\renewcommand{\cfttoctitlefont}{\Huge\bfseries}
\renewcommand{\cftbeforetoctitleskip}{2em}
\renewcommand{\cftaftertoctitleskip}{1.2em}
\renewcommand{\cftchapfont}{\large\bfseries}
\renewcommand{\cftchappagefont}{\large}
\renewcommand{\cftsecfont}{\normalsize}
\renewcommand{\cftsubsecfont}{\normalsize}
\renewcommand{\cftbeforechapskip}{1.2em}
\renewcommand{\cftbeforesecskip}{0.4em}

% ---------- 超链接 ----------
\usepackage[colorlinks=true,linkcolor=link,citecolor=link,urlcolor=link]{hyperref}

% ---------- 图形/表格题注 ----------
\usepackage{caption}
\captionsetup{font=small, labelfont=bf}

% ---------- 知识块（饱和色标题条 + 同色左边线 + 极淡底 + 方角） ----------
\usepackage[most]{tcolorbox}
\tcbset{
  jbox/.style = {
    enhanced, breakable, arc=3pt,
    boxrule=0.6pt, leftrule=3pt,   % 粗左边线（同标题色）
    left=10pt, right=10pt, top=8pt, bottom=8pt,
    before skip=12pt plus 2pt, after skip=12pt plus 2pt,
    fonttitle=\bfseries, coltitle=white,
    toptitle=1.5pt, bottomtitle=1.5pt,
    theorem style=standard,
    separator sign none, terminator sign none,
    description delimiters={\kern-0.42em（}{）},
    % parbox 用默认(true)：盒内文字自动按盒宽换行，不会溢出边框
  },
}
% 彩色定理环境：一行一个，颜色即第 3 个参数
%   用法：\newcoloredtheorem{环境名}{显示名}{颜色}{标签前缀}
\newcommand{\newcoloredtheorem}[4]{%
  \newtcbtheorem[number within=chapter]{#1}{#2}{jbox, colback=#3!8, colframe=#3, colbacktitle=#3}{#4}%
}
\newcoloredtheorem{definition}{定义}{bl}{def}
\newcoloredtheorem{theorem}{定理}{rd}{thm}
\newcoloredtheorem{lemma}{引理}{in}{lem}
\newcoloredtheorem{proposition}{命题}{pu}{prop}
\newcoloredtheorem{corollary}{推论}{pk}{cor}
\newcoloredtheorem{example}{例}{gr}{ex}
\newcoloredtheorem{remark}{注}{gy}{rem}
% 通用彩色信息块：#1=颜色(可选，默认蓝) #2=标题(必填)
%   用法：\begin{notebox}[颜色]{标题} ... \end{notebox}
\newtcolorbox{notebox}[2][bl]{jbox, colback=#1!8, colframe=#1, colbacktitle=#1, title={#2}}

% ---------- 代码框（语法高亮 + 代码框） ----------
% 语法高亮用 listings，外面套 tcolorbox 做个带标题的代码框。
%   用法：\begin{codebox}[语言]{标题}  代码  \end{codebox}
\usepackage{listings}
\definecolor{codebg}{HTML}{F5F6F8}   % 代码区背景
\definecolor{codeframe}{HTML}{9AA2B0} % 代码框边框
\definecolor{kw}{HTML}{4A6FD0}        % 关键字
\definecolor{st}{HTML}{4C9F70}        % 字符串
\definecolor{cm}{HTML}{9AA0A6}        % 注释
\definecolor{nm}{HTML}{E8933A}        % 行号
\lstset{
  basicstyle=\ttfamily\small,
  keywordstyle=\color{kw}\bfseries,
  stringstyle=\color{st},
  commentstyle=\color{cm}\itshape,
  numberstyle=\color{nm}\tiny,
  numbers=left, numbersep=6pt,
  showstringspaces=false, breaklines=true, tabsize=2,
  frame=none,
}
% listings 不带内置的 latex 语言，自定义一个；常见命令可高亮
\lstdefinelanguage{latex}{%
  morekeywords={documentclass,usepackage,begin,end,chapter,section,subsection,subsubsection,
    label,ref,eqref,caption,include,input,equation,align,figure,tabular,hline,centering,
    includegraphics,textbf,textit,emph,texttt,maketitle,tableofcontents,titlepage,body},
  morecomment=[l]{\%},
  sensitive=true,
}
\lstalias{tex}{latex}
\lstalias{sh}{bash}
% 无高亮的纯文本（目录/说明用），避免把 ​# 当注释导致解析问题
\lstdefinelanguage{text}{sensitive=true}
\newtcblisting{codebox}[2][bash]{%
  enhanced, breakable, listing only,
  listing options={language=#1},
  colback=codebg, colframe=codeframe, leftrule=3pt,
  colbacktitle=bl!80!black, coltitle=white, fonttitle=\bfseries,
  arc=2mm, boxrule=0.5pt,
  left=6pt, right=6pt, top=4pt, bottom=4pt,
  title={#2},
}
  引入（在 \begin{document} 之前）。
%   统一改风格（字体、配色、页面、框）都是改这里。
% ============================================================

% ---------- 字体（XeLaTeX，已由 ctex 加载 xeCJK/fontspec） ----------
% 正文思源宋体；中文「斜体」按国内惯例映射成楷体（\textit / \emph 即楷体）
\setCJKmainfont[ItalicFont=LXGW WenKai, BoldItalicFont=LXGW WenKai]{Noto Serif CJK SC}
\setCJKsansfont{Noto Sans CJK SC}    % 标题中文：思源黑体
\setCJKmonofont{LXGW WenKai Mono}    % 中文等宽：霞鹜文楷 Mono
\setmonofont{JetBrainsMono NF}       % 拉丁等宽（代码）：JetBrains Mono Nerd Font
% 拉丁正文与数学保留 Computer Modern。

% ---------- 数学 ----------
\usepackage{amsmath,amssymb}
\usepackage{bm}
\usepackage{mathtools}
\numberwithin{equation}{chapter}   % 公式编号按章：1.1 / 1.2 ...
\numberwithin{figure}{chapter}     % 图编号按章

% ---------- 页面（A4） ----------
\usepackage[a4paper, top=2.3cm, bottom=2.3cm, left=2.2cm, right=2.2cm,
  headheight=14pt, headsep=6mm, footskip=11mm]{geometry}
\usepackage{setspace}
\onehalfspacing

% ---------- 图形 ----------
\usepackage{graphicx}
\usepackage{tikz}
\usetikzlibrary{calc, arrows.meta, positioning, shapes.geometric}

% ---------- 表格 ----------
\usepackage{booktabs}
\usepackage{array}

% ---------- 物理单位 ----------
\usepackage{siunitx}

% ---------- 列表间距 ----------
\usepackage{enumitem}
\setlist{topsep=3pt, partopsep=0pt, itemsep=2pt, parsep=0pt}

% ---------- 颜色（知识块调色板 + 封面） ----------
\usepackage{xcolor}
\definecolor{bl}{HTML}{3B6BB0}   % 蓝 — 定义
\definecolor{tl}{HTML}{1F9B8E}   % 青 — 理解
\definecolor{gr}{HTML}{4C9F70}   % 绿 — 例
\definecolor{og}{HTML}{E8933A}   % 橙 — 证明思路
\definecolor{pu}{HTML}{7B5EA7}   % 紫 — 命题
\definecolor{rd}{HTML}{C0504D}   % 红 — 定理
\definecolor{in}{HTML}{283593}   % 靛 — 引理
\definecolor{pk}{HTML}{C86B98}   % 粉 — 推论
\definecolor{gy}{HTML}{6B7280}   % 灰 — 注
\definecolor{coverblue}{HTML}{283593}   % 封面深蓝
\definecolor{coverlight}{HTML}{5C6BC0}  % 封面浅蓝
\definecolor{covergray}{HTML}{5A5F6E}   % 封面灰
\definecolor{link}{HTML}{1F6F8B}

% ---------- 页眉页脚 ----------
\usepackage{fancyhdr}
\pagestyle{fancy}
\fancyhf{}
\fancyhead[L]{\small\nouppercase{\rightmark}}
\fancyhead[R]{\small\thepage}
\renewcommand{\headrulewidth}{0.4pt}

% ---------- 章节标题（加大醒目） ----------
\usepackage{titlesec}
\titleformat{\chapter}{\centering\fontsize{20}{26}\selectfont\bfseries}{\Huge 第\chinese{chapter}章}{0.6em}{}
\titlespacing*{\chapter}{0pt}{-16pt}{30pt}
\titleformat{\section}{\fontsize{16}{20}\selectfont\bfseries}{\thesection}{0.6em}{}
\titlespacing*{\section}{0pt}{16pt}{6pt}
\titleformat{\subsection}{\fontsize{14}{18}\selectfont\bfseries}{\thesubsection}{0.6em}{}
\titlespacing*{\subsection}{0pt}{12pt}{4pt}
\titleformat{\subsubsection}{\fontsize{12}{15}\selectfont\bfseries}{\thesubsubsection}{0.6em}{}

% ---------- 目录（加大字号） ----------
\usepackage{tocloft}
\renewcommand{\cfttoctitlefont}{\Huge\bfseries}
\renewcommand{\cftbeforetoctitleskip}{2em}
\renewcommand{\cftaftertoctitleskip}{1.2em}
\renewcommand{\cftchapfont}{\large\bfseries}
\renewcommand{\cftchappagefont}{\large}
\renewcommand{\cftsecfont}{\normalsize}
\renewcommand{\cftsubsecfont}{\normalsize}
\renewcommand{\cftbeforechapskip}{1.2em}
\renewcommand{\cftbeforesecskip}{0.4em}

% ---------- 超链接 ----------
\usepackage[colorlinks=true,linkcolor=link,citecolor=link,urlcolor=link]{hyperref}

% ---------- 图形/表格题注 ----------
\usepackage{caption}
\captionsetup{font=small, labelfont=bf}

% ---------- 知识块（饱和色标题条 + 同色左边线 + 极淡底 + 方角） ----------
\usepackage[most]{tcolorbox}
\tcbset{
  jbox/.style = {
    enhanced, breakable, arc=3pt,
    boxrule=0.6pt, leftrule=3pt,   % 粗左边线（同标题色）
    left=10pt, right=10pt, top=8pt, bottom=8pt,
    before skip=12pt plus 2pt, after skip=12pt plus 2pt,
    fonttitle=\bfseries, coltitle=white,
    toptitle=1.5pt, bottomtitle=1.5pt,
    theorem style=standard,
    separator sign none, terminator sign none,
    description delimiters={\kern-0.42em（}{）},
    % parbox 用默认(true)：盒内文字自动按盒宽换行，不会溢出边框
  },
}
% 彩色定理环境：一行一个，颜色即第 3 个参数
%   用法：\newcoloredtheorem{环境名}{显示名}{颜色}{标签前缀}
\newcommand{\newcoloredtheorem}[4]{%
  \newtcbtheorem[number within=chapter]{#1}{#2}{jbox, colback=#3!8, colframe=#3, colbacktitle=#3}{#4}%
}
\newcoloredtheorem{definition}{定义}{bl}{def}
\newcoloredtheorem{theorem}{定理}{rd}{thm}
\newcoloredtheorem{lemma}{引理}{in}{lem}
\newcoloredtheorem{proposition}{命题}{pu}{prop}
\newcoloredtheorem{corollary}{推论}{pk}{cor}
\newcoloredtheorem{example}{例}{gr}{ex}
\newcoloredtheorem{remark}{注}{gy}{rem}
% 通用彩色信息块：#1=颜色(可选，默认蓝) #2=标题(必填)
%   用法：\begin{notebox}[颜色]{标题} ... \end{notebox}
\newtcolorbox{notebox}[2][bl]{jbox, colback=#1!8, colframe=#1, colbacktitle=#1, title={#2}}

% ---------- 代码框（语法高亮 + 代码框） ----------
% 语法高亮用 listings，外面套 tcolorbox 做个带标题的代码框。
%   用法：\begin{codebox}[语言]{标题}  代码  \end{codebox}
\usepackage{listings}
\definecolor{codebg}{HTML}{F5F6F8}   % 代码区背景
\definecolor{codeframe}{HTML}{9AA2B0} % 代码框边框
\definecolor{kw}{HTML}{4A6FD0}        % 关键字
\definecolor{st}{HTML}{4C9F70}        % 字符串
\definecolor{cm}{HTML}{9AA0A6}        % 注释
\definecolor{nm}{HTML}{E8933A}        % 行号
\lstset{
  basicstyle=\ttfamily\small,
  keywordstyle=\color{kw}\bfseries,
  stringstyle=\color{st},
  commentstyle=\color{cm}\itshape,
  numberstyle=\color{nm}\tiny,
  numbers=left, numbersep=6pt,
  showstringspaces=false, breaklines=true, tabsize=2,
  frame=none,
}
% listings 不带内置的 latex 语言，自定义一个；常见命令可高亮
\lstdefinelanguage{latex}{%
  morekeywords={documentclass,usepackage,begin,end,chapter,section,subsection,subsubsection,
    label,ref,eqref,caption,include,input,equation,align,figure,tabular,hline,centering,
    includegraphics,textbf,textit,emph,texttt,maketitle,tableofcontents,titlepage,body},
  morecomment=[l]{\%},
  sensitive=true,
}
\lstalias{tex}{latex}
\lstalias{sh}{bash}
% 无高亮的纯文本（目录/说明用），避免把 ​# 当注释导致解析问题
\lstdefinelanguage{text}{sensitive=true}
\newtcblisting{codebox}[2][bash]{%
  enhanced, breakable, listing only,
  listing options={language=#1},
  colback=codebg, colframe=codeframe, leftrule=3pt,
  colbacktitle=bl!80!black, coltitle=white, fonttitle=\bfseries,
  arc=2mm, boxrule=0.5pt,
  left=6pt, right=6pt, top=4pt, bottom=4pt,
  title={#2},
}
  引入（在 \begin{document} 之前）。
%   统一改风格（字体、配色、页面、框）都是改这里。
% ============================================================

% ---------- 字体（XeLaTeX，已由 ctex 加载 xeCJK/fontspec） ----------
% 正文思源宋体；中文「斜体」按国内惯例映射成楷体（\textit / \emph 即楷体）
\setCJKmainfont[ItalicFont=LXGW WenKai, BoldItalicFont=LXGW WenKai]{Noto Serif CJK SC}
\setCJKsansfont{Noto Sans CJK SC}    % 标题中文：思源黑体
\setCJKmonofont{LXGW WenKai Mono}    % 中文等宽：霞鹜文楷 Mono
\setmonofont{JetBrainsMono NF}       % 拉丁等宽（代码）：JetBrains Mono Nerd Font
% 拉丁正文与数学保留 Computer Modern。

% ---------- 数学 ----------
\usepackage{amsmath,amssymb}
\usepackage{bm}
\usepackage{mathtools}
\numberwithin{equation}{chapter}   % 公式编号按章：1.1 / 1.2 ...
\numberwithin{figure}{chapter}     % 图编号按章

% ---------- 页面（A4） ----------
\usepackage[a4paper, top=2.3cm, bottom=2.3cm, left=2.2cm, right=2.2cm,
  headheight=14pt, headsep=6mm, footskip=11mm]{geometry}
\usepackage{setspace}
\onehalfspacing

% ---------- 图形 ----------
\usepackage{graphicx}
\usepackage{tikz}
\usetikzlibrary{calc, arrows.meta, positioning, shapes.geometric}

% ---------- 表格 ----------
\usepackage{booktabs}
\usepackage{array}

% ---------- 物理单位 ----------
\usepackage{siunitx}

% ---------- 列表间距 ----------
\usepackage{enumitem}
\setlist{topsep=3pt, partopsep=0pt, itemsep=2pt, parsep=0pt}

% ---------- 颜色（知识块调色板 + 封面） ----------
\usepackage{xcolor}
\definecolor{bl}{HTML}{3B6BB0}   % 蓝 — 定义
\definecolor{tl}{HTML}{1F9B8E}   % 青 — 理解
\definecolor{gr}{HTML}{4C9F70}   % 绿 — 例
\definecolor{og}{HTML}{E8933A}   % 橙 — 证明思路
\definecolor{pu}{HTML}{7B5EA7}   % 紫 — 命题
\definecolor{rd}{HTML}{C0504D}   % 红 — 定理
\definecolor{in}{HTML}{283593}   % 靛 — 引理
\definecolor{pk}{HTML}{C86B98}   % 粉 — 推论
\definecolor{gy}{HTML}{6B7280}   % 灰 — 注
\definecolor{coverblue}{HTML}{283593}   % 封面深蓝
\definecolor{coverlight}{HTML}{5C6BC0}  % 封面浅蓝
\definecolor{covergray}{HTML}{5A5F6E}   % 封面灰
\definecolor{link}{HTML}{1F6F8B}

% ---------- 页眉页脚 ----------
\usepackage{fancyhdr}
\pagestyle{fancy}
\fancyhf{}
\fancyhead[L]{\small\nouppercase{\rightmark}}
\fancyhead[R]{\small\thepage}
\renewcommand{\headrulewidth}{0.4pt}

% ---------- 章节标题（加大醒目） ----------
\usepackage{titlesec}
\titleformat{\chapter}{\centering\fontsize{20}{26}\selectfont\bfseries}{\Huge 第\chinese{chapter}章}{0.6em}{}
\titlespacing*{\chapter}{0pt}{-16pt}{30pt}
\titleformat{\section}{\fontsize{16}{20}\selectfont\bfseries}{\thesection}{0.6em}{}
\titlespacing*{\section}{0pt}{16pt}{6pt}
\titleformat{\subsection}{\fontsize{14}{18}\selectfont\bfseries}{\thesubsection}{0.6em}{}
\titlespacing*{\subsection}{0pt}{12pt}{4pt}
\titleformat{\subsubsection}{\fontsize{12}{15}\selectfont\bfseries}{\thesubsubsection}{0.6em}{}

% ---------- 目录（加大字号） ----------
\usepackage{tocloft}
\renewcommand{\cfttoctitlefont}{\Huge\bfseries}
\renewcommand{\cftbeforetoctitleskip}{2em}
\renewcommand{\cftaftertoctitleskip}{1.2em}
\renewcommand{\cftchapfont}{\large\bfseries}
\renewcommand{\cftchappagefont}{\large}
\renewcommand{\cftsecfont}{\normalsize}
\renewcommand{\cftsubsecfont}{\normalsize}
\renewcommand{\cftbeforechapskip}{1.2em}
\renewcommand{\cftbeforesecskip}{0.4em}

% ---------- 超链接 ----------
\usepackage[colorlinks=true,linkcolor=link,citecolor=link,urlcolor=link]{hyperref}

% ---------- 图形/表格题注 ----------
\usepackage{caption}
\captionsetup{font=small, labelfont=bf}

% ---------- 知识块（饱和色标题条 + 同色左边线 + 极淡底 + 方角） ----------
\usepackage[most]{tcolorbox}
\tcbset{
  jbox/.style = {
    enhanced, breakable, arc=3pt,
    boxrule=0.6pt, leftrule=3pt,   % 粗左边线（同标题色）
    left=10pt, right=10pt, top=8pt, bottom=8pt,
    before skip=12pt plus 2pt, after skip=12pt plus 2pt,
    fonttitle=\bfseries, coltitle=white,
    toptitle=1.5pt, bottomtitle=1.5pt,
    theorem style=standard,
    separator sign none, terminator sign none,
    description delimiters={\kern-0.42em（}{）},
    % parbox 用默认(true)：盒内文字自动按盒宽换行，不会溢出边框
  },
}
% 彩色定理环境：一行一个，颜色即第 3 个参数
%   用法：\newcoloredtheorem{环境名}{显示名}{颜色}{标签前缀}
\newcommand{\newcoloredtheorem}[4]{%
  \newtcbtheorem[number within=chapter]{#1}{#2}{jbox, colback=#3!8, colframe=#3, colbacktitle=#3}{#4}%
}
\newcoloredtheorem{definition}{定义}{bl}{def}
\newcoloredtheorem{theorem}{定理}{rd}{thm}
\newcoloredtheorem{lemma}{引理}{in}{lem}
\newcoloredtheorem{proposition}{命题}{pu}{prop}
\newcoloredtheorem{corollary}{推论}{pk}{cor}
\newcoloredtheorem{example}{例}{gr}{ex}
\newcoloredtheorem{remark}{注}{gy}{rem}
% 通用彩色信息块：#1=颜色(可选，默认蓝) #2=标题(必填)
%   用法：\begin{notebox}[颜色]{标题} ... \end{notebox}
\newtcolorbox{notebox}[2][bl]{jbox, colback=#1!8, colframe=#1, colbacktitle=#1, title={#2}}

% ---------- 代码框（语法高亮 + 代码框） ----------
% 语法高亮用 listings，外面套 tcolorbox 做个带标题的代码框。
%   用法：\begin{codebox}[语言]{标题}  代码  \end{codebox}
\usepackage{listings}
\definecolor{codebg}{HTML}{F5F6F8}   % 代码区背景
\definecolor{codeframe}{HTML}{9AA2B0} % 代码框边框
\definecolor{kw}{HTML}{4A6FD0}        % 关键字
\definecolor{st}{HTML}{4C9F70}        % 字符串
\definecolor{cm}{HTML}{9AA0A6}        % 注释
\definecolor{nm}{HTML}{E8933A}        % 行号
\lstset{
  basicstyle=\ttfamily\small,
  keywordstyle=\color{kw}\bfseries,
  stringstyle=\color{st},
  commentstyle=\color{cm}\itshape,
  numberstyle=\color{nm}\tiny,
  numbers=left, numbersep=6pt,
  showstringspaces=false, breaklines=true, tabsize=2,
  frame=none,
}
% listings 不带内置的 latex 语言，自定义一个；常见命令可高亮
\lstdefinelanguage{latex}{%
  morekeywords={documentclass,usepackage,begin,end,chapter,section,subsection,subsubsection,
    label,ref,eqref,caption,include,input,equation,align,figure,tabular,hline,centering,
    includegraphics,textbf,textit,emph,texttt,maketitle,tableofcontents,titlepage,body},
  morecomment=[l]{\%},
  sensitive=true,
}
\lstalias{tex}{latex}
\lstalias{sh}{bash}
% 无高亮的纯文本（目录/说明用），避免把 ​# 当注释导致解析问题
\lstdefinelanguage{text}{sensitive=true}
\newtcblisting{codebox}[2][bash]{%
  enhanced, breakable, listing only,
  listing options={language=#1},
  colback=codebg, colframe=codeframe, leftrule=3pt,
  colbacktitle=bl!80!black, coltitle=white, fonttitle=\bfseries,
  arc=2mm, boxrule=0.5pt,
  left=6pt, right=6pt, top=4pt, bottom=4pt,
  title={#2},
}
  引入（在 \begin{document} 之前）。
%   统一改风格（字体、配色、页面、框）都是改这里。
% ============================================================

% ---------- 字体（XeLaTeX，已由 ctex 加载 xeCJK/fontspec） ----------
% 正文思源宋体；中文「斜体」按国内惯例映射成楷体（\textit / \emph 即楷体）
\setCJKmainfont[ItalicFont=LXGW WenKai, BoldItalicFont=LXGW WenKai]{Noto Serif CJK SC}
\setCJKsansfont{Noto Sans CJK SC}    % 标题中文：思源黑体
\setCJKmonofont{LXGW WenKai Mono}    % 中文等宽：霞鹜文楷 Mono
\setmonofont{JetBrainsMono NF}       % 拉丁等宽（代码）：JetBrains Mono Nerd Font
% 拉丁正文与数学保留 Computer Modern。

% ---------- 数学 ----------
\usepackage{amsmath,amssymb}
\usepackage{bm}
\usepackage{mathtools}
\numberwithin{equation}{chapter}   % 公式编号按章：1.1 / 1.2 ...
\numberwithin{figure}{chapter}     % 图编号按章

% ---------- 页面（A4） ----------
\usepackage[a4paper, top=2.3cm, bottom=2.3cm, left=2.2cm, right=2.2cm,
  headheight=14pt, headsep=6mm, footskip=11mm]{geometry}
\usepackage{setspace}
\onehalfspacing

% ---------- 图形 ----------
\usepackage{graphicx}
\usepackage{tikz}
\usetikzlibrary{calc, arrows.meta, positioning, shapes.geometric}

% ---------- 表格 ----------
\usepackage{booktabs}
\usepackage{array}

% ---------- 物理单位 ----------
\usepackage{siunitx}

% ---------- 列表间距 ----------
\usepackage{enumitem}
\setlist{topsep=3pt, partopsep=0pt, itemsep=2pt, parsep=0pt}

% ---------- 颜色（知识块调色板 + 封面） ----------
\usepackage{xcolor}
\definecolor{bl}{HTML}{3B6BB0}   % 蓝 — 定义
\definecolor{tl}{HTML}{1F9B8E}   % 青 — 理解
\definecolor{gr}{HTML}{4C9F70}   % 绿 — 例
\definecolor{og}{HTML}{E8933A}   % 橙 — 证明思路
\definecolor{pu}{HTML}{7B5EA7}   % 紫 — 命题
\definecolor{rd}{HTML}{C0504D}   % 红 — 定理
\definecolor{in}{HTML}{283593}   % 靛 — 引理
\definecolor{pk}{HTML}{C86B98}   % 粉 — 推论
\definecolor{gy}{HTML}{6B7280}   % 灰 — 注
\definecolor{coverblue}{HTML}{283593}   % 封面深蓝
\definecolor{coverlight}{HTML}{5C6BC0}  % 封面浅蓝
\definecolor{covergray}{HTML}{5A5F6E}   % 封面灰
\definecolor{link}{HTML}{1F6F8B}

% ---------- 页眉页脚 ----------
\usepackage{fancyhdr}
\pagestyle{fancy}
\fancyhf{}
\fancyhead[L]{\small\nouppercase{\rightmark}}
\fancyhead[R]{\small\thepage}
\renewcommand{\headrulewidth}{0.4pt}

% ---------- 章节标题（加大醒目） ----------
\usepackage{titlesec}
\titleformat{\chapter}{\centering\fontsize{20}{26}\selectfont\bfseries}{\Huge 第\chinese{chapter}章}{0.6em}{}
\titlespacing*{\chapter}{0pt}{-16pt}{30pt}
\titleformat{\section}{\fontsize{16}{20}\selectfont\bfseries}{\thesection}{0.6em}{}
\titlespacing*{\section}{0pt}{16pt}{6pt}
\titleformat{\subsection}{\fontsize{14}{18}\selectfont\bfseries}{\thesubsection}{0.6em}{}
\titlespacing*{\subsection}{0pt}{12pt}{4pt}
\titleformat{\subsubsection}{\fontsize{12}{15}\selectfont\bfseries}{\thesubsubsection}{0.6em}{}

% ---------- 目录（加大字号） ----------
\usepackage{tocloft}
\renewcommand{\cfttoctitlefont}{\Huge\bfseries}
\renewcommand{\cftbeforetoctitleskip}{2em}
\renewcommand{\cftaftertoctitleskip}{1.2em}
\renewcommand{\cftchapfont}{\large\bfseries}
\renewcommand{\cftchappagefont}{\large}
\renewcommand{\cftsecfont}{\normalsize}
\renewcommand{\cftsubsecfont}{\normalsize}
\renewcommand{\cftbeforechapskip}{1.2em}
\renewcommand{\cftbeforesecskip}{0.4em}

% ---------- 超链接 ----------
\usepackage[colorlinks=true,linkcolor=link,citecolor=link,urlcolor=link]{hyperref}

% ---------- 图形/表格题注 ----------
\usepackage{caption}
\captionsetup{font=small, labelfont=bf}

% ---------- 知识块（饱和色标题条 + 同色左边线 + 极淡底 + 方角） ----------
\usepackage[most]{tcolorbox}
\tcbset{
  jbox/.style = {
    enhanced, breakable, arc=3pt,
    boxrule=0.6pt, leftrule=3pt,   % 粗左边线（同标题色）
    left=10pt, right=10pt, top=8pt, bottom=8pt,
    before skip=12pt plus 2pt, after skip=12pt plus 2pt,
    fonttitle=\bfseries, coltitle=white,
    toptitle=1.5pt, bottomtitle=1.5pt,
    theorem style=standard,
    separator sign none, terminator sign none,
    description delimiters={\kern-0.42em（}{）},
    % parbox 用默认(true)：盒内文字自动按盒宽换行，不会溢出边框
  },
}
% 彩色定理环境：一行一个，颜色即第 3 个参数
%   用法：\newcoloredtheorem{环境名}{显示名}{颜色}{标签前缀}
\newcommand{\newcoloredtheorem}[4]{%
  \newtcbtheorem[number within=chapter]{#1}{#2}{jbox, colback=#3!8, colframe=#3, colbacktitle=#3}{#4}%
}
\newcoloredtheorem{definition}{定义}{bl}{def}
\newcoloredtheorem{theorem}{定理}{rd}{thm}
\newcoloredtheorem{lemma}{引理}{in}{lem}
\newcoloredtheorem{proposition}{命题}{pu}{prop}
\newcoloredtheorem{corollary}{推论}{pk}{cor}
\newcoloredtheorem{example}{例}{gr}{ex}
\newcoloredtheorem{remark}{注}{gy}{rem}
% 通用彩色信息块：#1=颜色(可选，默认蓝) #2=标题(必填)
%   用法：\begin{notebox}[颜色]{标题} ... \end{notebox}
\newtcolorbox{notebox}[2][bl]{jbox, colback=#1!8, colframe=#1, colbacktitle=#1, title={#2}}

% ---------- 代码框（语法高亮 + 代码框） ----------
% 语法高亮用 listings，外面套 tcolorbox 做个带标题的代码框。
%   用法：\begin{codebox}[语言]{标题}  代码  \end{codebox}
\usepackage{listings}
\definecolor{codebg}{HTML}{F5F6F8}   % 代码区背景
\definecolor{codeframe}{HTML}{9AA2B0} % 代码框边框
\definecolor{kw}{HTML}{4A6FD0}        % 关键字
\definecolor{st}{HTML}{4C9F70}        % 字符串
\definecolor{cm}{HTML}{9AA0A6}        % 注释
\definecolor{nm}{HTML}{E8933A}        % 行号
\lstset{
  basicstyle=\ttfamily\small,
  keywordstyle=\color{kw}\bfseries,
  stringstyle=\color{st},
  commentstyle=\color{cm}\itshape,
  numberstyle=\color{nm}\tiny,
  numbers=left, numbersep=6pt,
  showstringspaces=false, breaklines=true, tabsize=2,
  frame=none,
}
% listings 不带内置的 latex 语言，自定义一个；常见命令可高亮
\lstdefinelanguage{latex}{%
  morekeywords={documentclass,usepackage,begin,end,chapter,section,subsection,subsubsection,
    label,ref,eqref,caption,include,input,equation,align,figure,tabular,hline,centering,
    includegraphics,textbf,textit,emph,texttt,maketitle,tableofcontents,titlepage,body},
  morecomment=[l]{\%},
  sensitive=true,
}
\lstalias{tex}{latex}
\lstalias{sh}{bash}
% 无高亮的纯文本（目录/说明用），避免把 ​# 当注释导致解析问题
\lstdefinelanguage{text}{sensitive=true}
\newtcblisting{codebox}[2][bash]{%
  enhanced, breakable, listing only,
  listing options={language=#1},
  colback=codebg, colframe=codeframe, leftrule=3pt,
  colbacktitle=bl!80!black, coltitle=white, fonttitle=\bfseries,
  arc=2mm, boxrule=0.5pt,
  left=6pt, right=6pt, top=4pt, bottom=4pt,
  title={#2},
}
